\section{Case Study: Qwen3-30B-A3B on DGX Spark}\label{sec:case2}

The second case is a mixture-of-experts model on the same machine, chosen because its decode step
looks nothing like the first: 48 layers, 128 experts of which 8 are routed per token, block FP8
weights ($128\times128$ scales), vLLM 0.34.0, no speculation. Here the books found no waste to
remove; what they did find is where the time goes, and that led to a different kind of fix.
Data and scripts are in \texttt{case-studies/gb10-qwen3-30b-a3b}.

\paragraph{Nothing is written for nothing.}
The access tracer, with weights skipped, finds that a batch-1 eager decode step writes 47.0\,MB of
non-weight memory, of which 0.4\,MB is overwritten before anyone reads it; at batch 4, 190.8 and
1.5\,MB. What little is dead comes from the kernel that lays tokens out per expert for the grouped
GEMM. Traced with CUDA graphs on (1\,169 launches per step, all seen, 96\% with parameters) the step
writes 78.9\,MB, 0.3\,MB of it dead. In the graph trace the DeepGEMM kernels, which move their
operands through tensor maps and encode their activation maps for 4\,096 rows, are taken to read up
to the end of the allocation segment each map starts in, which hides the scatter kernel's dead
bytes; the eager trace is the more precise one here.

\paragraph{The experts are read exactly, but slowly.}
Because the expert weights move through tensor maps, the tracer cannot say which experts a step
reads. An Nsight Compute profile of the same eager step can say how much each launch pulled from
memory into L2 (\cref{tab:moe-traffic}). Each layer's grouped GEMMs pull 2--6\% more than the eight
routed experts' weights, so routing costs little extra traffic. They run at 60--66\% of
the ceiling, while the dense FP8 GEMMs of the first case reach 85--91\%; the bytes are right and
the kernel is slow. The step pulls 3.49\,GB, 13.8\,ms at the ceiling, against 30\,ms of GPU time
in eager mode; the expert GEMMs lose about 4\,ms of it and \texttt{lm\_head} 1.6\,ms.

\begin{table}[t]
\centering\small
\begin{tabular}{@{}lrrrrr@{}}
\toprule
Launch & per layer & pulled & needed & time & of ceiling \\
\midrule
Grouped GEMM, gate and up & 1 & 26.8\,MB & 25.2\,MB & 160\,\textmu s & 66\% \\
Grouped GEMM, down & 1 & 12.8\,MB & 12.6\,MB & 84\,\textmu s & 60\% \\
\texttt{lm\_head} ($151936\times2048$ BF16) & 1 per step & 633\,MB & 622\,MB & 4.1\,ms & 61\% \\
\bottomrule
\end{tabular}
\caption{Bytes pulled into L2 (Nsight Compute, \texttt{lts\_\_d\_sectors\_fill\_sysmem}) against the weight bytes of the eight routed experts, one eager decode step at batch 1.}\label{tab:moe-traffic}
\end{table}

\paragraph{Other backends.}
vLLM chooses the MoE backend from a priority list. On Hopper it moves Triton to the front for
block FP8 without expert parallelism; SM12x has no such rule and falls through to DeepGEMM.
\Cref{tab:moe-backend} runs the same model with each backend, CUDA graphs on. Triton keeps the
W8A8 path and is 12.5\% faster per decode step at batch 1 and 19.3\% at batch 4. Marlin is as fast at batch 1 and faster at batch 4, but dequantizes the weights and computes in BF16, which changes numerics.

\begin{table}[t]
\centering\small
\begin{tabular}{@{}lrr@{}}
\toprule
\texttt{--moe-backend} & decode step, batch 1 & decode step, batch 4 \\
\midrule
\texttt{deep\_gemm} (default) & 23.52\,ms & 48.02\,ms \\
\texttt{triton} & 20.58\,ms ($-12.5\%$) & 38.73\,ms ($-19.3\%$) \\
\texttt{marlin} (weight-only FP8) & 20.74\,ms ($-11.8\%$) & 36.34\,ms ($-24.3\%$) \\
\texttt{cutlass} & \multicolumn{2}{c}{not supported for this configuration} \\
\bottomrule
\end{tabular}
\caption{Qwen3-30B-A3B FP8 decode step by MoE backend, vLLM 0.34.0, CUDA graphs.}\label{tab:moe-backend}
\end{table}

\paragraph{Serving.}
With \texttt{vllm bench serve} (random 512 input and 256 output tokens, one run per point) Triton
raises throughput from 45.32 to 54.97\,tok/s at concurrency 1 ($+21.3\%$), from 99.63 to
117.61 at concurrency 4 ($+18.0\%$) and from 136.90 to 196.55 at concurrency 8 ($+43.6\%$). In a
prefill-bound throughput run (4\,096 input, 16 output tokens) the two differ by 1\%. GSM8K
(200 questions, 5-shot, greedy) is 0.965 with DeepGEMM and 0.960 with Triton. We proposed
preferring Triton for block FP8 MoE without expert parallelism (tensor parallelism,
including a single GPU) on SM12x, as on Hopper~\cite{vllmpr69355}. Only GB10 was measured; other SM12x parts share the kernels but not the
memory system.

\paragraph{What the case shows.}
In the first case the books found bytes that should not have been written; in this one they found
none, and showed instead that every byte read was needed and the reading was slow. Both answers
come from the same accounting: per-launch bytes, the ceiling, and, where needed, what each launch
read, wrote and pulled.
